% !TeX program = lualatex
% Compile twice: lualatex open_problems_500.tex
% All references and prose are embedded; no BibTeX database or external inputs.
\documentclass[11pt,a4paper]{article}
\usepackage[margin=24mm,headheight=14pt]{geometry}
\usepackage{fontspec}
\setmainfont{Latin Modern Roman}
\setsansfont{Latin Modern Sans}
\setmonofont{Latin Modern Mono}
\usepackage{amsmath,amssymb,mathtools}
\usepackage{unicode-math}
\setmathfont{Latin Modern Math}
\usepackage{microtype}
\usepackage{enumitem,xurl,needspace}
\usepackage[dvipsnames]{xcolor}
\usepackage{hyperref}
\hypersetup{unicode=true,colorlinks=true,linkcolor=MidnightBlue,urlcolor=MidnightBlue,
 pdftitle={Research notebook: Problems 40--49},pdfauthor={Source-annotated compilation},bookmarksnumbered=true}
\usepackage{bookmark}
\usepackage{tocloft}
\setlength{\cftsecnumwidth}{3em}
\cftsetpnumwidth{2.5em}
\cftsetrmarg{4em}
\usepackage{titlesec}
\titleformat{\section}{\Large\bfseries\raggedright}{\thesection}{0.7em}{}
\usepackage{fancyhdr}
\pagestyle{fancy}\fancyhf{}
\fancyhead[L]{\small\sffamily Mathematical problem statements}
\fancyhead[R]{\small Source edition 22}
\fancyfoot[C]{\thepage}
\setlength{\parindent}{0pt}
\setlength{\parskip}{5pt plus 1pt}
\setlength{\emergencystretch}{3em}
\setcounter{tocdepth}{1}
\setcounter{secnumdepth}{1}
\setlist[itemize]{leftmargin=3em,itemsep=5pt,topsep=4pt,parsep=0pt}
\allowdisplaybreaks
\newcommand{\N}{\mathbb{N}}
\newcommand{\Z}{\mathbb{Z}}
\newcommand{\Q}{\mathbb{Q}}
\newcommand{\R}{\mathbb{R}}
\newcommand{\C}{\mathbb{C}}
\newcommand{\F}{\mathbb{F}}
\newcommand{\A}{\mathbb{A}}
\newcommand{\PP}{\mathbb P}
\newcommand{\E}{\mathbb{E}}
\newcommand{\eps}{\varepsilon}
\newcommand{\dd}{\,\mathrm d}
\newcommand{\sref}[2]{\hyperlink{source:#1:#2}{[S#2]}}
\newcommand{\eref}[2]{\hyperlink{extra:#1:#2}{[E#2]}}
% Turn the next switch off for a shorter reading copy. The quotations preserve
% every exactTarget field, including qualifications omitted from a short summary.
\newif\ifshowcatalogue
\showcataloguetrue
\newcommand{\cataloguescope}[1]{\ifshowcatalogue
 \paragraph{Exact scope in the supplied catalogue.}
 \begin{quote}\small #1\end{quote}\fi}

% Research additions dated 27 September 2026; compile twice with LuaLaTeX.
\numberwithin{equation}{section}
\begin{document}
\setcounter{section}{41}
% ================================================================
% problemId: problem.conformal-field-theory-limit-of-the-critical-three-dimensional-ising-model
% source releaseRank: 42; source releaseStatus: open
% exactTarget SHA-256: 71c2e81c6b7353dd4e8b9caafb0b196abeacf7ad424c9f2b32820e802c588f7a
\clearpage
\section{\texorpdfstring{Conformal Field Theory Limit of the Critical Three-Dimensional Ising Model}{Conformal Field Theory Limit of the Critical Three-Dimensional Ising Model}}\label{42}
\subsection{Definitions and mathematical statement}
For Ising spins $\sigma_x\in\{-1,1\}$ on a three-dimensional lattice, the zero-field Gibbs weight is proportional to $\exp(\beta\sum_{\langle x,y\rangle}\sigma_x\sigma_y)$. At critical inverse temperature $\beta_c$, rescale the lattice spacing to $a\downarrow0$ and normalize the spin field as a random distribution, schematically
\[
 \Phi_a(f)=c_a\sum_{x\in a\Z^3}f(x)\sigma_x.
\]
The goal is a nontrivial limiting field with conformally covariant correlation functions and non-Gaussian statistics. For example, a centered Gaussian limit would satisfy $\langle\Phi_1\Phi_2\Phi_3\Phi_4\rangle=\sum_{\rm pairings}\langle\Phi_i\Phi_j\rangle\langle\Phi_k\Phi_l\rangle$. The selected target requires failure of such Wick factorization. Normalization and convergence topology are part of a rigorous scaling-limit construction, not specified numerically by the catalogue.
\par\smallskip\textit{Formulation and concept references:} \sref{42}{1}.
\cataloguescope{Prove that the scaling limit of the critical three-dimensional Ising model exists and is conformally invariant, with a non-Gaussian spin field whose limiting spin correlations do not factorize according to Wick's rule.}
\subsection{Short English statement}
Does the critical three-dimensional Ising model converge to a genuinely interacting conformally invariant field theory?
% BEGIN RESEARCH ADDITION 42
\subsection{Research attempt}
\paragraph{Current frontier.}
The requested conclusion contains three distinct tasks: construction of a
scaling limit, conformal covariance of that limit, and survival of a connected
four-point function. The survey in \sref{42}{1}, Sections 6.3 and 8.4, explains
both the random-current description of non-Gaussianity and the outstanding
three-dimensional conformal-limit problem. The two-point estimates of
Duminil-Copin--Panis imply divergence of the critical bubble diagram in
three and four dimensions; their bound $\eta\le 1/2$ in dimension three is
conditional on existence of that critical exponent, not a construction of
it \sref{42}{3}, \eref{42}{1}. Bootstrap calculations of a three-dimensional
four-point function are informative within their conformal-field-theory
framework but do not prove convergence of the lattice model
\eref{42}{2}. These inputs are kept separate below.

\paragraph{Attack route.}
One route is to normalize by an actual lattice variance rather than assume
an unproved critical exponent. This makes compactness relatively accessible,
but leaves uniqueness and the scale dependence of the normalization.
A second route is to use switching of random currents to quantify exactly
what must persist for a non-Gaussian limit. A third, genuinely additional,
step would establish conformal covariance; neither compactness nor a
nonzero cumulant supplies that step. We pursue the first two routes and
make the missing probabilistic estimate explicit.

\paragraph{Attempt: a normalization and a subsequential construction.}
Work in a translation-invariant, spin-flip-symmetric, zero-field infinite-volume
critical Ising state obtained by the usual ferromagnetic limiting procedure.
Write $G(u,v)=\langle\sigma_u\sigma_v\rangle\ge0$ on $\Z^3$.
The positivity and four-point inequality used here are the standard
ferromagnetic inequalities and the switching identity reviewed in
\sref{42}{1}, Sections 6.1--6.3. The following compactness argument is an
auxiliary deduction; it is not claimed as a historically new scaling-limit
theorem.

For $0<a\le1/2$, set $L=\lfloor a^{-1}\rfloor$ and
\[
 B_L=\{0,\ldots,L-1\}^3,\quad M_L=\sum_{u\in B_L}\sigma_u,
 \quad b_L^2=\langle M_L^2\rangle,\quad
 \Phi_a(f)=b_L^{-1}\sum_{u\in\Z^3}f(au)\sigma_u.
\]
Since all covariances are nonnegative and $\sigma_u^2=1$, $b_L^2\ge L^3$.
In particular the normalization is defined and uses no assumed power law.
Tile $\Z^3$ by translates of $B_L$. A fixed compact physical set $K$
intersects at most $N_K$ such blocks, uniformly in small $a$, because
$aL\ge1/2$. Translation invariance gives $L^2$ norm one to every normalized
block sum. For any possibly complex $f$ supported in $K$, positivity of $G$
and then the triangle inequality in $L^2$ give
\begin{align*}
 \langle |\Phi_a(f)|^2\rangle
 &\le \|f\|_\infty^2 b_L^{-2}
       \left\langle\left(\sum_{u\text{ in the covering blocks}}
                      \sigma_u\right)^2\right\rangle\\
 &\le N_K^2\|f\|_\infty^2. \tag{42.1}
\end{align*}
There is no assumption of independence between blocks in this estimate.

Here is a topology in which (42.1) gives tightness. Put a smooth cutoff
$\chi$ inside a fixed larger box, identify that box with a torus, and let
$e_k$ be its Fourier basis. Its functions have uniformly bounded sup norm.
For $s>3/2$, (42.1) and the definition of the negative Sobolev norm imply
\[
 \sup_a\E\|\chi\Phi_a\|_{H^{-s}}^2
 \le C_\chi\sum_{k\in\Z^3}(1+|k|^2)^{-s}<\infty. \tag{42.2}
\]
Fix $s'>s$. Bounded sets in $H^{-s}$ are relatively compact in $H^{-s'}$:
the high-frequency tail of their squared $H^{-s'}$ norm is bounded by
$(1+R^2)^{-(s'-s)}$ times their squared $H^{-s}$ norm, and the remaining
Fourier coordinates form a finite-dimensional bounded set. Markov's
inequality applied to (42.2) therefore gives compact containment. Using a
countable family of cutoffs exhausting $\R^3$, with failure probabilities
summing to an arbitrarily small number, gives tightness in
$H^{-s'}_{\rm loc}(\R^3)$. Consequently every sequence $a_j\downarrow0$
has a subsequence converging in law in this space. This proves only a
subsequential statement with the displayed normalization.

The subsequential limits are not identically zero. Choose smooth,
nonnegative, compactly supported $f_0$ with $f_0\ge1$ on $[0,1]^3$.
Covariance positivity gives $\E\Phi_a(f_0)^2\ge1$. The four-spin inequality
implies, for nonnegative weights,
\[
 \E\Phi_a(f_0)^4\le3\bigl(\E\Phi_a(f_0)^2\bigr)^2\le C_{f_0}.
 \tag{42.3}
\]
For distinct sites this inequality follows from the nonpositive connected
four-point function. Coincident indices cause no problem: for example,
$\langle\sigma_i^2\sigma_j\sigma_k\rangle=G(j,k)$ is at most
$G(j,k)+2G(i,j)G(i,k)$. Summing the inequalities proves (42.3).
Thus the squares are uniformly integrable. Evaluation at $f_0$ is
continuous in the chosen local distribution topology, so passage along
any convergent subsequence preserves the lower bound on its second
moment. This establishes nontriviality, not non-Gaussianity.

\paragraph{Attempt: an exact non-Gaussianity certificate.}
For four distinct sites, let $U_4$ denote the four-point function minus its
three Wick pairings. The switching identity in \sref{42}{1}, equation
(6.20), is
\[
 \begin{aligned}
 U_4(u_1,u_2,u_3,u_4)
 &=-2G(u_1,u_2)G(u_3,u_4)\\
 &\quad\times
 \bigl(\mathbb P^{\{u_1,u_2\}}\!\otimes
   \mathbb P^{\{u_3,u_4\}}\bigr)[\mathcal C],
 \end{aligned}
 \tag{42.4}
\]
where $\mathcal C$ is the event that all four sites are connected in
$\mathbf n_1+\mathbf n_2$. The probabilities here are source-conditioned random-current probabilities,
not independent random-walk probabilities. One may first use the finite-volume
identity and pass to the chosen state; equivalently the bounded connection
factor is the corresponding limiting ratio where the denominator is positive.

Choose four nonnegative smooth test functions with pairwise separated
compact supports. Put $X_i^{(a)}=\Phi_a(f_i)$ and
$B_{ij}^{(a)}=\E X_i^{(a)}X_j^{(a)}$. When
$B_{12}^{(a)}B_{34}^{(a)}>0$, define $p_a$ to be the weighted average of the
connection probabilities in (42.4), with weights
\[
 b_L^{-4}\prod_{i=1}^4 f_i(au_i)G(u_1,u_2)G(u_3,u_4),
\]
divided by $B_{12}^{(a)}B_{34}^{(a)}$. All sums are finite, all weights are
nonnegative, and their sum equals that denominator. Therefore $0\le p_a\le1$
and summing (42.4) proves the exact identity
\[
 \kappa_4(X_1^{(a)},X_2^{(a)},X_3^{(a)},X_4^{(a)})
       =-2B_{12}^{(a)}B_{34}^{(a)}p_a. \tag{42.5}
\]
This isolates a useful conditional criterion. Suppose, along a convergent
subsequence, that $B_{12}^{(a)}\to b_{12}>0$,
$B_{34}^{(a)}\to b_{34}>0$, and $\liminf p_a\ge p_*>0$.
Assume also uniform eighth moments for these four smeared fields. This last
assumption is stated explicitly here rather than inferred from (42.3).
H\"older's inequality then gives uniform integrability of their mixed
fourth products. Joint convergence and moment convergence in (42.5) imply
\[
 \kappa_4(X_1,X_2,X_3,X_4)\le-2b_{12}b_{34}p_*<0.
 \tag{42.6}
\]
The limiting field consequently fails Wick factorization. The outstanding
substantive input for this route is a uniform lower bound for the
\emph{weighted, source-conditioned} connection probability, together with
nondegeneracy of the separated covariances. The unit-block normalization
alone establishes neither.

\paragraph{Stress test.}
Finite-lattice non-Gaussianity does not survive automatically. For $N$
independent symmetric spins, $Y_N=N^{-1/2}\sum_{i=1}^N\sigma_i$ has
$\E Y_N^2=1$ and
\[
 \E Y_N^4=3-2/N,\qquad \kappa_4(Y_N)=-2/N\longrightarrow0.
\]
This exact calculation is checked by enumeration for $N\le8$ in the
accompanying script. It is a counterexample to a proposed inference, not
to the interacting critical Ising statement. Similarly, divergence of a
bubble sum is not itself a lower bound for $p_a$: the cited bubble result
also covers dimension four, where Gaussian scaling behavior is established
for the relevant Ising limits \sref{42}{1}, Section 6.3; \sref{42}{3}.
Source conditioning and the normalization in (42.5) cannot be dropped in a
random-walk intersection analogy.

The compactness argument deliberately leaves major information undetermined.
Different subsequences may give different fields; $b_L$ has not been shown
to have regular variation; the limit has not been shown to be rotation,
dilation, or inversion covariant. Even translation covariance inherited in
a limit would not establish the full conformal group. Neither a chosen
normalization nor a nonzero cumulant supplies a conformal transformation
law. These omissions are part of the target, not optional refinements.

\paragraph{Outcome.}
\textbf{Outcome: Exploratory approach with unresolved gap.}
A direct variance-normalized tightness argument gives nontrivial
subsequential random-distribution limits, and an exact random-current
identity reduces persistence of a specified non-Gaussian four-point
function to a positive weighted connection probability. The conditional
criterion and the elementary compactness proof are not assertions of
historical novelty. Uniqueness across scales, the required connection
estimate, and conformal covariance remain unproved here. Thus the full
scaling-limit and conformal-invariance assertion has not been established.
% END RESEARCH ADDITION 42
\subsection{Sources}
\begin{itemize}\raggedright
\item[\textnormal{[S1]}]\hypertarget{source:42:1}{} 100 Years of the (Critical) Ising Model on the Hypercubic Lattice (Hugo Duminil-Copin).
\url{https://arxiv.org/html/2208.00864}.
{\small\textit{Catalogue formulation source.}}
\item[\textnormal{[S2]}]\hypertarget{source:42:2}{} The Ising model: highlights and perspectives.
\url{https://link.springer.com/article/10.1007/s11040-025-09515-1}.
{\small\textit{Catalogue research/status reference; not a proof endorsement.}}
\item[\textnormal{[S3]}]\hypertarget{source:42:3}{} New lower bounds for the (near) critical Ising and phi4 models' two-point functions.
\url{https://www.unige.ch/~duminil/publi/2024_lower_bound_Ising.pdf}.
{\small\textit{Catalogue research/status reference; not a proof endorsement.}}
\item[\textnormal{[S4]}]\hypertarget{source:42:4}{} Certified reduced-basis emulation of conformal field theories on the fuzzy sphere.
\url{https://arxiv.org/abs/2609.11675}.
{\small\textit{Catalogue research/status reference; not a proof endorsement.}}
% BEGIN RESEARCH SOURCES 42
\item[\textnormal{[E1]}]\hypertarget{extra:42:1}{} Hugo Duminil-Copin and Romain Panis,
\emph{New lower bounds for the (near) critical Ising and $\varphi^4$ models' two-point functions},
arXiv:2404.05700v1 (2024), Theorems 1.5 and 1.8; author PDF dated 9 April 2024.
\url{https://arxiv.org/abs/2404.05700}.
{\small\textit{Identification of the manuscript retained as [S3]; exponent qualification and bubble divergence. Consulted 27 September 2026.}}
\item[\textnormal{[E2]}]\hypertarget{extra:42:2}{} Slava Rychkov, David Simmons-Duffin and Bernardo Zan,
\emph{Non-gaussianity of the critical 3d Ising model}, arXiv:1612.02436 (2016), abstract and four-point-function analysis.
\url{https://arxiv.org/abs/1612.02436}.
{\small\textit{Conformal-bootstrap analysis, not an unconditional lattice scaling-limit construction. Consulted 27 September 2026.}}
% END RESEARCH SOURCES 42
\end{itemize}

\end{document}
